\origpage{550}{N. Robertson, P. Seymour / Journal of Combinatorial Theory, Series B 102 (2012) 530--563}
\noindent replaced by
\begin{equation*}
p, \theta', \Gamma, K_1, \Sigma, \mathcal{T}, d, Z', F_1, \ldots, F_t, R, \Sigma', \Gamma_3, L_2 \cup (L \cap K_1)
\end{equation*}
respectively. We recall that $\theta'$ was chosen so that 5.3 holds with $p, \theta$ replaced by $p, \theta'$. To verify the hypotheses of 5.3 is straightforward. 5.3(i) follows from 7.3(v); 5.3(ii) from 7.3(iv); 5.3(iii) from (2); and the other hypotheses follow from (6) and (7). Consequently, by 5.3, there exists $L'_0 \subseteq (\Gamma \cap \Sigma') \cup K_1$ with the same effect on $Z'$ as $L_2 \cup (L \cap K_1)$ and with $L'_0 \cap K_1 \subseteq L_2 \cup (L \cap K_1)$. Now $L_2 \cap K_1 = L \cap K_1$, and so $L'_0 \cap K_1 \subseteq L$. Let $f \in \phi(E(H)) \cap E(L'_0)$. Since $\phi(E(H)) \cap E(L) = \emptyset$, it follows that $f \notin E(L)$, and hence $f \notin E(K_1)$. By (1), $f \in E(A_1 \cup \cdots \cup A_t)$. Consequently,
\begin{equation*}
\phi(E(H)) \cap E(L'_0) \subseteq E(A_1 \cup \cdots \cup A_t).
\end{equation*}
This proves (8).

Let $Z = Z_0 \cup \bar{\pi}(A_1) \cup \cdots \cup \bar{\pi}(A_t)$, and define $\alpha \colon Z' \to Z$ as follows: if $v \in Z_0$, let $\alpha(v) = v$, and if $v \in \bar{\pi}(A'_i)$ where $1 \le i \le t$, and $v$ is the $j$th term of $\pi(A'_i)$ say where $1 \le j \le |\bar{\pi}_i|$, let $\alpha(v)$ be the $j$th term of $\pi(A_i)$. This defines a function since the sets $Z_0, \bar{\pi}(A'_1), \ldots, \bar{\pi}(A'_t)$ are mutually disjoint. Similarly, for $v \in Z'$ define $\mu(v) = v$ if $v \in Z_0$, and if $v$ is the $j$th term of $\pi(A'_i)$ let $\mu(v)$ be the $j$th term of $\pi_i$. Thus, if $v \in \bar{\pi}(A'_i)$, $\beta(\mu(v)) = v$.

\hspace*{12pt}(9) \textit{$u, v \in Z'$ are $L'_0$-connected if and only if $\alpha(u), \alpha(v)$ are $L_0 \cup (L \cap K_1)$-connected.}

\textit{Subproof.} To show this we make a sequence of equivalent statements, starting with:

\hspace*{12pt}(a) $\alpha(u), \alpha(v)$ are $L_0 \cup (L \cap K_1)$-connected.

Since $\alpha(u), \alpha(v) \in Z \subseteq V(L_0)$, (a) is equivalent to

\hspace*{12pt}(b) $\alpha(u), \alpha(v)$ are $L_1 \cup (L \cap K_1)$-connected,

because $L_1 \cup (L \cap K_1)$ is obtained from $L_0 \cup (L \cap K_1)$ by adding vertices of degree 1. Now $\alpha(u)$ and $\mu(u)$ are either equal or are adjacent in $L_1$; and similarly for $\alpha(v), \mu(v)$. Consequently, (b) is equivalent to

\hspace*{12pt}(c) $\mu(u), \mu(v)$ are $L_1 \cup (L \cap K_1)$-connected.

There is an isomorphism between $L_1 \cup (L \cap K_1)$ and $L_2 \cup (L \cap K_1)$ (since $\beta$ fixes $U(L_1 \cap (L \cap K_1))$ pointwise), mapping each vertex $x$ to $\beta(x)$ if $x \in \Sigma$ and mapping $x$ to itself otherwise. Since $\beta(\mu(v)) = v$ for $v \in \bar{\pi}(A'_1) \cup \cdots \cup \bar{\pi}(A'_t)$, this isomorphism maps $\mu(u)$ to $u$ and $\mu(v)$ to $v$. Consequently, (c) is equivalent to

\hspace*{12pt}(d) $u, v$ are $L_2 \cup (L \cap K_1)$-connected.

But by (8), (d) is equivalent to

\hspace*{12pt}(e) $u, v$ are $L'_0$-connected.

Hence (a) is equivalent to (e). This proves (9).

Let $A_{t+1}$ be the subdigraph of $G$ with vertex set $Z_0$ and edge set $\phi(E(H)) \cap E(K_1)$, let $\pi(A_{t+1})$ be some ordering of $Z_0 \setminus W$ and let $\bar{\pi}(A_{t+1}) = Z_0 \setminus W$. Let $A_0 = (\Gamma_0 \cup K_1) \setminus (\phi(E(H)) \cap E(K_1))$.

\hspace*{12pt}(10) \textit{$\{A_0, A_1, \ldots, A_{t+1}\}$ is a $Z$-division of $\Gamma \cup K_1 \cup A_1 \cup \cdots \cup A_t$, and $Z \subseteq V(A_0)$, and for $1 \le i \le t+1$, $\pi(A_i) + \omega$ is an ordering of $Z \cap V(A_i)$.}
